\chapter{The Problem of Disciplinary Dispersion}

\epigraph{\emph{What converges in the world disperses in the disciplines.}}{}

If the previous chapter identifies a recurring tendency, the present chapter names the obstacle that has prevented that tendency from becoming an explicit field of inquiry. The obstacle is not simply lack of data. It is disciplinary dispersion. The relevant findings are scattered across neuroscience, psychoacoustics, dynamical systems, ecology, biosemiotics, field theory, wave physics, and computational modeling, but the institutions of knowledge production remain organized around local objects, local methods, and local vocabularies. As a result, neighboring discoveries often remain mutually invisible even when they concern formally adjacent patterns of coupling, recurrence, and resonant organization.

The first barrier is lexical. A neuroscientist speaks of cross-frequency coupling, entrainment, or traveling waves. A psychoacoustician speaks of consonance, roughness, periodicity, or harmonicity. A physicist may instead speak of normal modes, field excitations, interference, coherence, or radiation. A complexity researcher speaks of recurrence, synchronization, and attractor structure. An ecologist may speak of niche partitioning, signal occupancy, or soundscape organization. These vocabularies are not merely superficial labels. They carry with them local ontologies, favored instruments, canonical datasets, and implicit criteria of relevance. What one field treats as a central explanatory variable, another may treat as a derivative artifact or ignore altogether. The result is not outright contradiction so much as non-recognition.

The second barrier is methodological localism. Each domain has built its own evidentiary apparatus around different scales and different objects. Neural recording emphasizes coupling statistics and anatomical interpretation. Psychoacoustics privileges controlled perceptual judgments and spectral models. Dynamical systems research focuses on formal stability and recurrence regimes. Wave physics may privilege interferometry, diffraction, spectral decomposition, or detection of propagating disturbances in media and fields. Ecology tracks population-level organization in real environments. These methodological commitments are justified within each field, but they also make cross-domain synthesis difficult. Evidence that is decisive in one literature may not even count as admissible evidence in another. A recurrence plot, an fMRI coupling map, an optical interference pattern, a gravitational-wave trace, and a perceptual consonance curve can all point toward structured relational organization while remaining institutionally non-convertible.

The third barrier is incentive structure. Scientific communities are rewarded for novelty within recognizable disciplinary boundaries, not for patiently assembling bridges across them. Journals, review cultures, citation practices, graduate training, and conference ecologies all reinforce local coherence. The researcher who frames a result in the accepted language of a subfield is legible, citable, and publishable. The researcher who attempts a transdisciplinary synthesis often risks being judged too speculative for one audience and too technical for another. The same formal pattern may be publishable when framed as cross-frequency coupling in neuroscience, legible when framed as recurrence in nonlinear dynamics, and suspect when framed as harmonicity across domains. In that environment, the safest intellectual strategy is partial rediscovery: solve a local problem with local vocabulary, and leave the larger pattern unnamed.

The result is that even disputes internal to one literature rarely propagate cleanly to adjacent ones. Critical debates on consonance have already shown how difficult it is to separate physiological, cultural, and formal explanatory layers without collapsing them into one another \parencite{parncutt_2011}. Cross-cultural work adds a second constraint. Findings such as those reported for the Tsimane do not invalidate the possibility of broad harmonic regularities, but they do block any easy slide from recurrent structure to universal aesthetic law \parencite{mcdermott_2016}. Similar caution is needed when one moves toward the hard sciences: the fact that multiple domains deploy wave and field concepts does not by itself license the claim that one model is structurally isomorphic to reality in any simple or exhaustive sense.

Galison names one part of this situation through trading zones: scientific subcultures can coordinate locally around exchanges, instruments, and partial contact languages without ever sharing a full ontology \parencite{galison_1997}. Star and Griesemer make the point more concrete. Some objects travel across communities precisely because they are plastic enough to be interpreted differently while remaining robust enough to stay recognizable \parencite{star_1989}. Harmonic relations, recurrence structures, constrained oscillatory couplings, and wave-like mode organizations may already function in practice as proto-boundary objects. The problem is that they have not yet been stabilized as such.

Lotman's semiosphere deepens the picture. Boundaries are not only lines of separation; they are also sites of translation, filtering, and productive asymmetry between heterogeneous spaces \parencite{lotman_2005}. If that is right, then disciplinary borders do not merely block circulation. They shape the terms under which circulation becomes possible. The problem for HIT is therefore not only the absence of contact. It is the absence of a sufficiently shared translational space in which adjacent findings can become mutually legible without losing their specificity.

Calvo's critique of structuralist-model isomorphism marks the limit that such a program must respect from the outset. A useful representational structure is not thereby a literal duplication of the world \parencite{calvo_2006}. HIT does not need to say that an acoustic ratio, a cortical eigenmode, and a field excitation are simply the same thing under different names. It needs the more exact claim that these domains may exhibit family resemblances, formal analogies, and partially translatable organizations worth coordinating without being collapsed into a single structural identity.

The epistemic cost of dispersion is high. First, fields repeatedly reinvent one another's partial insights under different names. Second, claims that should be compared remain incomparable because they are embedded in different measurement traditions. Third, overstatement and understatement coexist. One community may localize a phenomenon so tightly that its broader significance disappears from view. Another may generalize too quickly from a narrow context because it lacks contact with adjacent literatures that would impose constraints. The result is a fragmented landscape in which a real convergence can be simultaneously visible in pieces and absent as an articulated object.

HIT should therefore not be introduced as a finished theory that simply waits to be accepted. That would reproduce the very problem it seeks to solve by imposing premature unity on a heterogeneous empirical terrain. At this stage, HIT is better understood as an epistemic category and a research program: it names the possibility that a family of related findings across several domains can be coordinated into a rigorous inquiry about harmonic relations as informational, dynamical, and biological constraints. What it contributes first is not closure, but a shared problem-space.

Lakatos's language of research programmes gives that posture a precise structure. A program acquires shape when a hard core of commitments is named and a protective belt of revisable auxiliary hypotheses is allowed to vary around it \parencite{lakatos_1978}. In the present case, the hard core can be stated minimally: structured harmonic and resonant relations appear repeatedly as privileged constraints in physical, biological, perceptual, and computational organization, and that recurrence is scientifically worth investigating as such. Around that core remains a protective belt that is deliberately open: which mechanisms matter most, in which domains, under which scales, for which forms of information processing, and with what limits.

Kuhn's account of paradigms clarifies the status of the enterprise from another angle. HIT does not yet constitute a paradigm in any strong sense. It is a pre-paradigmatic convergence across several mature fields, and shared objects of inquiry do not automatically create a common research space \parencite{kuhn_1962}. That is why the program must be built neither as an imperial synthesis that seeks to subsume every neighboring language nor as a timid pluralism that refuses integration altogether.

Popper's demarcation problem adds the complementary requirement. If HIT is to become more than a suggestive umbrella, it must remain vulnerable to informative failure rather than insulated by analogy alone \parencite{popper_1959}. The framework therefore has to distinguish observation from mechanism, convergence from proof, analogy from isomorphism, and suggestive correspondence from experimentally supported inference. Part~II begins from that demand for conceptual precision. The next question is no longer whether a distributed convergence exists, but what, exactly, would count as harmonic information.
